\documentclass[UTF8,12pt,a4paper,fontset=fandol]{ctexart}

% 支持从项目根目录、深度学习目录或当前目录编译。
\makeatletter
\def\input@path{{../}{../../}}
\makeatother
\usepackage{researchnotes}

\title{模型量化的原理与前置知识}
\author{}
\date{}

\begin{document}
\maketitle

\begin{abstract}
模型量化用有限的数值集合表示神经网络中的权重、激活和缓存，在存储、访存、计算成本与预测质量之间进行取舍。理解量化，需要把数值表示、量化误差、网络敏感性、优化目标和硬件执行连接起来，而不能只把它看作浮点数到整数的类型转换。本文从可手算的例子出发，介绍二进制与浮点数、线性代数、概率统计、微分和优化等前置知识，推导均匀与非均匀量化、整数矩阵乘法、误差传播、校准、训练后量化和量化感知训练的基本原理，进一步解释二阶补偿、通道缩放、正交旋转、低秩适配与键值缓存量化。最后给出完整数值算例、实现流程、评估设计和练习解答。文中严格区分精确恒等式、依赖假设的近似和需要实验验证的工程判断；示例用于说明机制，不构成任何模型或设备上的性能报告。
\end{abstract}
\tableofcontents

\section{研究对象、核心问题与学习路线}
\label{sec:quant-intro}

\subsection{从四个权重开始理解量化}
假设一个线性层的四个权重为 $(-0.83,-0.21,0.36,0.91)$。选择间距 $s=0.1$，把每个权重除以 $s$ 并舍入到最近整数，得到 $(-8,-2,4,9)$。以后只保存这些整数和公共尺度 $s$，便可重建近似权重 $(-0.8,-0.2,0.4,0.9)$。四个重建误差依次为 $(0.03,0.01,0.04,-0.01)$。如果输入是 $(1,2,-1,1)^\top$，原来的内积为 $-0.70$，重建后的内积为 $-0.70$；权重发生了变化，某个输入上的输出却恰好不变，因为四项误差相互抵消。若输入换成 $(1,0,0,0)^\top$，输出误差就变成 $0.03$。这个例子说明，参数误差与函数误差不是同一个量，量化的好坏必须联系实际输入判断。

\keyterm{模型量化（model quantization）}是用有限的可表示数值集合近似模型中原有数值的过程。它可以作用于\keyterm{权重（weight）}、前向计算产生的\keyterm{激活（activation）}、注意力机制保存的缓存，也可以作用于训练中的梯度和优化器状态。量化方案至少要回答四个问题：允许哪些重建值，哪些元素共享尺度，如何选择重建值，以及执行时如何完成运算。只说明“使用四位模型”，并没有给出这些问题的完整答案。

量化主要利用数值近似来换取资源节省，通常保持网络主体的层数和连接结构。\keyterm{剪枝（pruning）}删除参数或结构，\keyterm{知识蒸馏（knowledge distillation）}改变学生模型的训练监督，\keyterm{低秩分解（low-rank factorization）}改变矩阵的因子表示，它们与量化可以组合，但不能互相替代。例如，稀疏矩阵如果仍按稠密格式保存，零元素依然占据空间；量化整数如果被装进三十二位容器而没有位打包，也不会自动得到理想的存储压缩。

\subsection{必须分开的几组概念}
\keyterm{位宽（bit width）}表示每个编码使用多少个二进制位，\keyterm{精度（precision）}描述可分辨数值的细致程度，\keyterm{动态范围（dynamic range）}描述可表示数值的跨度。相同位宽可以对应不同的范围和间距；浮点数与均匀整数编码更不是相同的数轴划分方式。存储精度、乘法输入精度和累加精度也应分别说明。低位宽输入经过较高精度累加，是常见而合理的计算组织，并不意味着量化失效。

本文使用 W4A16 表示权重的名义位宽为四、激活的名义位宽为十六；W8A8 同理。这种记号不指定整数还是浮点格式，也不指定累加器、偏置和缓存的格式。一个 W4A16 实现可以在读取权重后把它解码为半精度浮点数，再完成乘法；也可以使用支持相应混合格式的专用内核。必须检查实际执行路径，才能判断节省的是内存、带宽还是算术时间。

\keyterm{训练后量化（post-training quantization, PTQ）}从已经训练好的模型出发，通过统计、搜索或局部重建确定量化结果；\keyterm{量化感知训练（quantization-aware training, QAT）}在训练或微调中模拟量化，使参数适应离散约束。PTQ 不等于完全不使用数据，也不等于绝不进行数值优化；QAT 不等于所有训练状态都使用整数。两类方法的界线在具体文献中可能有细微差别，阅读时应直接核对哪些参数被更新、用了什么数据和损失。

\subsection{全文路线与记号约定}
前半部分先回答“数如何表示”和“误差如何计算”，再把标量量化推广到矩阵与网络；后半部分围绕如何降低损失和如何得到真实部署收益展开。只关心基础实现的读者，可以先掌握第\ref{sec:quant-number}、\ref{sec:quant-scalar}、\ref{sec:quant-granularity}、\ref{sec:quant-integer}节，再阅读完整算例；希望理解大模型论文的读者，应继续掌握重建目标、二阶补偿和等价变换。

除特别说明外，单个输入为列向量 $\mathbf x\in\mathbb R^d$，权重矩阵为 $\mathbf W\in\mathbb R^{m\times d}$，偏置为 $\boldsymbol\beta\in\mathbb R^m$，线性层输出为 $\mathbf y=\mathbf W\mathbf x+\boldsymbol\beta$。校准输入按列组成 $\mathbf X\in\mathbb R^{d\times N}$，其中 $N$ 是输入向量数；语言模型中一个词元位置通常就贡献一个输入向量。帽号表示量化后的重建值，如 $\widehat{\mathbf W}$；整数编码记为 $q$；量化误差一律定义为“重建值减原值”。位宽记为 $b$，尺度记为 $s>0$，整数零点记为 $z$。

\section{二进制、定点数与浮点数}
\label{sec:quant-number}

\subsection{整数编码和定点解释}
一个二进制位只有零和一两种状态，因此 $b$ 位一共有 $2^b$ 种编码。无符号整数通常解释为 $0,1,\ldots,2^b-1$；二进制补码通常解释为 $-2^{b-1},\ldots,2^{b-1}-1$。因此八位有符号整数的通常范围是 $[-128,127]$，并非 $[-127,127]$。量化算法可以主动不用其中一个编码，以获得对称数值集合，但必须说明这种约定，不能把算法选择误当成整数类型本身的定义。

\keyterm{定点数（fixed-point number）}可以理解为带有固定尺度的整数。例如整数 $q=37$ 与尺度 $2^{-4}$ 共同表示 $37/16$。若尺度限定为二的整数次幂，乘尺度可能通过移位及相应定点处理实现；一般的神经网络量化尺度不必是二的幂。仿射量化采用 $s(q-z)$ 的重建规则，其中 $z$ 把某个整数编码对应到实数零。整数编码是离散标签，重建规则才赋予它数值意义，所以不能仅根据编码大小比较来自不同尺度的两个张量。

对同一个整数类型，改变尺度会同时改变分辨率和可覆盖范围。以对称八位编码 $q\in\{-127,\ldots,127\}$ 为例，$s=0.01$ 时范围是 $[-1.27,1.27]$，相邻重建值差 $0.01$；$s=0.1$ 时范围扩大十倍，间距也扩大十倍。量化最基本的困难已经出现：既想保留很大的异常值，又想精细刻画密集的小数值，在固定编码数量下通常不能同时做到。

\subsection{浮点数的有效数字与指数}
一个二进制正规浮点数可写成 $(-1)^\sigma(1.f)_2\,2^e$，其中 $\sigma$ 是符号位，$f$ 是尾数的小数部分，$e$ 是扣除偏置后的指数。指数使数轴间距随数值大小变化：在相同指数区间内等间距，在更大的指数区间间距随之放大。\keyterm{次正规数（subnormal number）}用于靠近零的区域，不再使用同样的隐含首位一规则。无穷大、非数值和溢出处理也属于格式定义的一部分，不能只根据指数位数和尾数位数推断所有特殊值。

常见 FP32 有一位符号、八位指数和二十三位显式尾数；FP16 有一位符号、五位指数和十位显式尾数；BF16 有一位符号、八位指数和七位显式尾数。FP16 比 BF16 提供更多有效数字位，BF16 则具有更宽的指数范围。因此把 BF16 一概叫作“比 FP16 更精确”是不准确的：对于相近数量级的两个数，FP16 可能更容易区分；对于很大或很小的数，BF16 可能更容易避免溢出或下溢。实际结果还取决于设备如何处理次正规数和舍入。

在舍入到最近值且没有溢出、下溢等异常的正规数范围内，浮点运算常用模型 $\operatorname{fl}(x)=x(1+\delta)$、$|\delta|\leq u$ 描述，其中 $u$ 是与有效位数有关的单位舍入误差。这个模型主要控制相对误差。均匀量化在不饱和时控制的则是绝对误差 $s/2$。一个很小但非零的数被舍入到零，相对误差可以达到百分之百，而绝对误差仍很小；因此比较格式时必须先说明关注哪种误差。

\subsection{低位浮点数与共享尺度}
FP8 是一类八位浮点格式的统称。文献\cite{micikevicius2022}讨论的 E4M3 和 E5M2 分别分配四位、五位指数及三位、两位显式尾数，连同符号位构成八位编码；其中特殊值约定也有差异。不能把所有名为 E4M3 的变体都当作完全相同的编码规范，更不能只凭“八位”推断它与 INT8 的误差一致。浮点格式更适合描述跨数量级的相对分辨率，整数仿射量化则显式用尺度适配局部范围。

低位浮点量化也经常引入外部尺度，先将张量除以尺度，再转换到低位浮点格式。\keyterm{块浮点（block floating point）}及共享尺度格式进一步让一个小块共用指数或尺度，以较少的元数据调整局部动态范围。块越小，通常越能适配局部数值分布，但尺度数量和处理开销增加。所谓格式选择，实际上同时涉及元素编码、共享尺度、块形状和算术支持，而不是仅选择一个数据类型名称。

\subsection{舍入、溢出与数值可复现性}
最近舍入还要规定恰好落在两个重建值中点时如何处理，例如取偶数、远离零或固定方向。取偶数可以避免某些连续同向舍入偏差，但不保证神经网络误差无偏。负数的右移、截断到零和向负无穷取整也不同，例如 $-3/2$ 截断到零得到 $-1$，向负无穷取整得到 $-2$。这些细节在整数重定标中会影响逐元素结果，参考实现与部署内核必须采用一致规则。

\keyterm{饱和（saturation）}是把越界值限制在可表示端点；整数溢出的环绕行为则可能把很大的正数变为负数。量化器通常需要前者而不是后者。判断代码是否正确，不能只测试靠近零的温和输入，还要覆盖端点、中点、负数、全零张量和极端累加长度。本文例题若未特别指出中点情况，均选择没有舍入歧义的数值。

\section{线性代数、概率与优化的必要基础}
\label{sec:quant-math}

\subsection{内积、范数和矩阵方向}
向量内积 $\mathbf a^\top\mathbf b=\sum_i a_i b_i$ 把每个坐标的乘积相加，因此不同坐标上的误差可以抵消，也可以叠加。欧氏范数定义为 $\|\mathbf a\|_2=(\sum_i a_i^2)^{1/2}$，无穷范数为 $\|\mathbf a\|_\infty=\max_i|a_i|$；矩阵的 Frobenius 范数为 $\|\mathbf A\|_F=(\sum_{ij}A_{ij}^2)^{1/2}$。这些范数分别刻画总体能量、最大坐标和矩阵逐元素误差，回答的问题不同。

矩阵的谱范数定义为 $\|\mathbf A\|_2=\sup_{\|\mathbf x\|_2=1}\|\mathbf A\mathbf x\|_2$，描述最坏输入方向上的放大倍数。由定义，对任意 $\mathbf x$ 都有 $\|\mathbf A\mathbf x\|_2\leq\|\mathbf A\|_2\|\mathbf x\|_2$，且 $\|\mathbf A\|_2\leq\|\mathbf A\|_F$。因此很小的逐元素误差在高维累加后未必仍然很小；另一方面，最坏方向如果在实际数据中几乎不出现，谱范数界又可能过于保守。确定性界与数据分布上的平均误差应同时理解。

量化中的“按通道”依赖矩阵方向。按本文 $\mathbf W\in\mathbb R^{m\times d}$ 的约定，一行对应一个输出通道，一列对应一个输入通道。若某个框架用 $\mathbf X^\top\mathbf W^\top$ 组织批量乘法，存储轴可能改变，数学含义并没有改变。实现时必须把轴编号翻译成输入或输出特征，而不能直接照抄另一份代码的维度数字。

\subsection{期望、方差与二阶矩}
设随机输入向量为 $\mathbf X_{\rm r}$，小写 $\mathbf x$ 表示其实际取值。均值为 $\boldsymbol\mu=\mathbb E\mathbf X_{\rm r}$，协方差为 $\boldsymbol\Sigma=\mathbb E[(\mathbf X_{\rm r}-\boldsymbol\mu)(\mathbf X_{\rm r}-\boldsymbol\mu)^\top]$。假设各坐标具有有限二阶矩，则未中心化二阶矩为
\begin{equation}
 \mathbf M=\mathbb E[\mathbf X_{\rm r}\mathbf X_{\rm r}^\top]
 =\boldsymbol\Sigma+\boldsymbol\mu\boldsymbol\mu^\top.
 \label{eq:quant-second-moment}
\end{equation}
除非均值为零，二阶矩不能直接叫作协方差。后文的线性层重建使用 $\mathbf M$；漏掉均值项可能改变最优补偿方向，尤其当激活明显非中心化时。

对固定权重误差 $\boldsymbol\delta$，线性输出误差是 $E=\boldsymbol\delta^\top\mathbf X_{\rm r}$，其均方值满足
\begin{equation}
 \mathbb E E^2=\boldsymbol\delta^\top\mathbf M\boldsymbol\delta.
 \label{eq:quant-output-mse}
\end{equation}
这是展开有限求和并应用期望线性性得到的精确等式，不要求输入为高斯分布，也不要求坐标独立。若 $\mathbf M$ 的非对角项很大，两个权重误差的联合影响就不能由各自平方误差单独描述。量化算法利用校准输入，正是为了估计网络实际关心的方向。

样本二阶矩 $\widehat{\mathbf M}=\mathbf X\mathbf X^\top/N$ 是对总体量的估计。独立同分布和有限二阶矩是常见的一致性分析条件，而实际语言序列中的相邻词元通常相关。即使估计在校准分布上准确，部署输入分布改变后也不保证有效。校准集的价值来自覆盖和代表性，并非仅来自样本数量。

\subsection{梯度、Hessian 与局部近似}
设网络参数组成向量 $\theta$，任务损失为 $L(\theta)$。梯度 $\nabla L$ 描述一阶变化，Hessian 矩阵 $\mathbf H=\nabla^2L$ 描述二阶变化。若 $L$ 在考察点邻域内二次连续可微，则当 $\|\Delta\theta\|_2\to0$ 时，
\begin{equation}
 L(\theta+\Delta\theta)-L(\theta)
 =\nabla L(\theta)^\top\Delta\theta
 +\tfrac12\Delta\theta^\top\mathbf H\Delta\theta
 +o(\|\Delta\theta\|_2^2).
 \label{eq:quant-taylor}
\end{equation}
只有在梯度足够小且扰动足够小的情况下，才适合忽略一阶项并用二阶项解释损失。训练结束不意味着所有梯度精确为零，四位或两位量化也未必属于小扰动，因而二阶分析应被视为有条件的近似工具。

对称矩阵 $\mathbf H$ 若对任意非零 $\mathbf v$ 满足 $\mathbf v^\top\mathbf H\mathbf v>0$，称为正定；若允许等于零，称为半正定。神经网络任务损失的 Hessian 不一定正定，而线性最小二乘目标的 Hessian 天然半正定。二者不能混淆。后文添加正数乘单位矩阵的\keyterm{阻尼（damping）}可保证局部二次问题可逆，同时也改变了所优化的代价。

\subsection{损失函数、训练与推理}
分类模型输出逻辑值向量 $\mathbf a\in\mathbb R^K$，经 $p_k=\exp(a_k)/\sum_j\exp(a_j)$ 得到 softmax 概率。对正确类别 $y$，交叉熵为 $-\log p_y$。回归中常用平方损失，语言模型训练则通常累加真实下一个词元的负对数概率。量化可以用这些任务损失直接训练，也可以只用中间层输出的平方误差作为替代目标；替代目标更便宜，但与最终任务质量并不等价。

训练时反向传播依据链式法则传播梯度，优化器还保存动量或二阶统计状态；推理时只执行前向计算。权重存储减少四倍，不能据此推断完整训练内存也减少四倍，因为激活保存、梯度、优化器状态和临时空间可能占据较大比例。理解 QAT 或量化微调前，必须先列清哪些对象需要长期保存、哪些对象只是临时用于运算。

\section{标量量化器：网格、零点与裁剪}
\label{sec:quant-scalar}

\subsection{仿射均匀量化的严格定义}
\begin{definition}[仿射均匀量化]
给定整数端点 $q_{\min}<q_{\max}$、尺度 $s>0$ 和整数零点 $z\in[q_{\min},q_{\max}]$，令 $\operatorname{clip}(u,a,c)=\min\{\max\{u,a\},c\}$，$\operatorname{round}$ 表示带有固定中点规则的最近舍入。编码器与解码器定义为
\begin{equation}
 q(x)=\operatorname{clip}\!\left(\operatorname{round}(x/s)+z,q_{\min},q_{\max}\right),
 \qquad \widehat x=Q(x)=s\bigl(q(x)-z\bigr).
 \label{eq:quant-affine}
\end{equation}
重建集合为 $\mathcal C=\{s(k-z):k=q_{\min},\ldots,q_{\max}\}$。
\end{definition}
这里的编码结果是整数，解码结果是实数近似。“反量化”通常就是这一解码操作，它不能恢复被舍入丢掉的信息。由于 $z$ 是可用整数，$Q(0)=0$，零被精确表示。这对卷积填充、稀疏激活和掩码处理有实际意义。没有整数字面零点约束的一般仿射码本也可以构造，但它与这里定义的标准整数零点方案不同。

给定希望覆盖的观测区间 $[a,c]$，常先令 $a_0=\min(a,0)$、$c_0=\max(c,0)$，再取
\begin{equation}
 s=\frac{c_0-a_0}{q_{\max}-q_{\min}},\qquad
 z=\operatorname{clip}\!\left(\operatorname{round}(q_{\min}-a_0/s),q_{\min},q_{\max}\right).
 \label{eq:quant-minmax}
\end{equation}
这里要求 $c_0>a_0$。由于 $z$ 需要舍入，实际重建端点是 $s(q_{\min}-z)$ 与 $s(q_{\max}-z)$，不一定恰好等于 $a_0,c_0$。若要求严格覆盖全部观测值，可在确定零点后相应调整尺度；不同库采用的端点调整规则需要单独确认。全零张量不能使用零尺度，应约定一个正尺度并把编码统一设为零点。

\subsection{对称量化和非对称量化}
\keyterm{对称量化（symmetric quantization）}通常围绕实数零设置对称重建网格。为避免端点不对称，本文在 $b\geq2$ 时采用 $q_{\max}=2^{b-1}-1$、$q_{\min}=-q_{\max}$、$z=0$，以最大绝对值 $\alpha>0$ 选 $s=\alpha/q_{\max}$。因此有符号 $b$ 位编码中有一个码字不用。也有实现使用完整补码范围，其正负端点并不完全对称；两种选择都可以实现，但量化阈值和误差计算必须与实际约定一致。

\keyterm{非对称量化（asymmetric quantization）}允许重建范围关于零不对称，适合明显偏移或单侧分布。例如非负激活取 $z=0$、无符号编码区间 $[0,255]$，便能将全部编码用于非负轴。这里“非对称”是指实数范围和网格，不表示输入概率分布一定偏斜，也不应与某些论文中键和值使用不同量化轴的“非对称设计”混为一谈。

\begin{example}[带零点的四位量化]
用无符号四位编码 $0,\ldots,15$ 表示区间 $[-1,2]$，取 $s=3/15=0.2$、$z=5$。对于 $x=0.74$，有 $q=9$、$\widehat x=0.8$，误差为 $0.06$；对于 $x=-0.93$，有 $q=0$、$\widehat x=-1$，误差为 $-0.07$；对于 $x=2.36$，编码饱和到 $15$，重建值为 $2$，误差为 $-0.36$。前两项误差不超过半个间距，最后一项因为越界而不受该界约束。
\end{example}

\subsection{不饱和误差界及其适用条件}
\begin{proposition}[最近网格投影的误差界]
设 $Q$ 如式\eqref{eq:quant-affine}，实际重建端点为 $\ell=s(q_{\min}-z)$ 和 $u=s(q_{\max}-z)$。对任意 $x\in[\ell,u]$，有 $|Q(x)-x|\leq s/2$；对任意实数 $x$，有
\begin{equation}
 |Q(x)-x|\leq \operatorname{dist}(x,[\ell,u])+s/2,
 \label{eq:quant-pointwise-bound}
\end{equation}
其中 $\operatorname{dist}(x,[\ell,u])$ 表示 $x$ 到区间的最短距离。
\end{proposition}
\begin{proof}
区间内相邻重建值的间距为 $s$，每个区间内的点距至少一个端点不超过 $s/2$，最近舍入选取这样的重建值。对任意 $x$，先投影到区间得到 $x_c$。三角不等式给出 $|Q(x)-x|\leq|Q(x)-x_c|+|x_c-x|$；区间外时 $Q(x)=x_c$，区间内时第一项至多为 $s/2$，得到所述界。
\end{proof}
该结论是绝对误差界，不能直接推出模型准确率变化的界。它也没有断言所有误差在区间内均匀分布。给定确定权重和确定尺度，每个误差都是确定数；只有引入输入分布、随机舍入或统计近似以后，才有理由谈误差的概率分布。

\subsection{随机舍入与无偏性的边界}
对位于相邻重建值 $ks$ 和 $(k+1)s$ 之间的 $x$，写 $x/s=k+r$，其中 $r\in[0,1]$。\keyterm{随机舍入（stochastic rounding）}以概率 $1-r$ 选择 $ks$，以概率 $r$ 选择 $(k+1)s$。在没有端点饱和时，条件期望满足 $\mathbb E[Q(x)\mid x]=x$，条件方差为 $s^2r(1-r)\leq s^2/4$。这是对舍入随机性的结论，不要求 $x$ 自身随机。

无偏不代表每次误差较小，也不代表更低的任务损失。对于非线性函数 $f$，一般有 $\mathbb E[f(Q(x))]\ne f(x)$；即使某一层无偏，经过下一层激活后也可能产生偏差。多个坐标使用相关随机数时，内积误差还会出现协方差项。随机舍入可以帮助训练中很小的更新不被长期截断，但需要考虑随机数生成成本和推理可复现性，不能只凭“无偏”判断它优于最近舍入。

\section{量化误差的统计模型与非均匀码本}
\label{sec:quant-distortion}

\subsection{均匀噪声近似从何而来}
令 $E=Q(X)-X$。在量化间距足够小、输入密度在每个网格内变化缓慢且饱和可忽略时，常把 $E$ 近似为 $[-s/2,s/2]$ 上的均匀随机变量，于是
\begin{equation}
 \mathbb E E\approx0,\qquad
 \mathbb E E^2\approx\frac1s\int_{-s/2}^{s/2}e^2\,\mathrm de=\frac{s^2}{12}.
 \label{eq:quant-uniform-noise}
\end{equation}
积分是该假设下的精确计算，但假设本身通常只是高分辨率近似。对高度离散、集中在网格附近、包含大量零值或只有两三位表示的张量，这个模型可能很不准确。噪声与信号独立也不是由均匀边缘分布自动推出的额外结论。

若保持重建区间长度近似不变，并且有效网格数量近似随 $2^b$ 增长，则 $s$ 近似正比于 $2^{-b}$，均方误差近似正比于 $2^{-2b}$。因此每增加一位，理想化均方误差约降到四分之一。用 $10\log_{10}(\mathbb E X^2/\mathbb E E^2)$ 定义信号量化噪声比，在信号功率固定等条件下每位约增加 $10\log_{10}4\approx6.02$ 分贝。这个数不是所有神经网络量化误差的普遍规律，端点裁剪和分布形状都会破坏其前提。

\subsection{裁剪为何可能降低总误差}
考虑对称实际端点 $\pm\alpha$、内部间距 $s=\alpha/q_{\max}$，并设随机变量 $X$ 具有密度 $p$ 和有限二阶矩。内部舍入误差与外部饱和误差分别贡献
\begin{equation}
 D(\alpha)=\int_{-\alpha}^{\alpha}(Q(x)-x)^2p(x)\,\mathrm dx
 +\int_{|x|>\alpha}(|x|-\alpha)^2p(x)\,\mathrm dx.
 \label{eq:quant-clipping-risk}
\end{equation}
这是以实际重建端点定义的精确分解。增大 $\alpha$ 通常降低尾部饱和误差，却扩大内部间距；减小 $\alpha$ 则反过来。最优范围并不一定覆盖样本极值，也不一定对应某个固定百分位。应当说明优化的是元素重建误差、层输出误差还是任务损失，因为它们可能选择不同阈值。

例如某个组中绝大多数数值的绝对值小于一，只有一个值为二十。若用对称四位量化并覆盖全部范围，间距接近 $20/7$，大量普通数值会映射到零。若裁剪到较小范围，这些普通数值的误差可以明显下降，但异常值会产生较大误差。异常值是否可以牺牲，取决于它出现的频率、所在通道的作用和后续网络敏感性；数值稀少并不等于任务贡献不重要。

\subsection{非均匀标量量化与质心条件}
均匀量化的重建值间距相同。\keyterm{非均匀量化（nonuniform quantization）}允许使用任意有序码本 $c_1<\cdots<c_K$，通过编码记录所选码字的索引。若固定码本并以平方误差为目标，最优分配选择最近码字，相邻决策边界为两码字中点。若固定每个分配区域 $R_k$，则最优重建值是条件均值：
\begin{equation}
 c_k=\frac{\int_{R_k}x p(x)\,\mathrm dx}{\int_{R_k}p(x)\,\mathrm dx},
 \label{eq:quant-centroid}
\end{equation}
条件是区域概率为正且相关矩存在。证明只需对区域内平方误差关于 $c_k$ 求导，令导数为零，再利用二阶导数为正。

交替更新最近邻区域与区域质心，便得到经典标量量化的迭代思路，也对应一维聚类中的平方误差最小化。每一步不增加当前失真，但一般只能保证到达局部稳定点，不能保证全局最优。空区域还需要单独处理。码本学习把更多重建值放到高密度区域，可能降低元素均方误差；代价是查表、码本元数据和更复杂的运算支持。

对数量级变化较大的数值，可以用对数量化，让相邻正码字具有近似固定比值。它更接近控制相对误差，但零无法直接取对数，需要专门编码，负数也需要符号处理。\keyterm{向量量化（vector quantization）}则把多个权重组成向量后共享一个码字索引，能够利用坐标相关性；其码本大小、检索成本、分组布局和解码方式都会影响真实收益。降低数学失真与提高硬件效率是两个同时需要满足的目标。

\subsection{率失真观点与有效位宽}
\keyterm{率失真（rate--distortion）}观点把表示成本和近似损失放在同一个优化问题中。对编码方案 $\mathcal Q$，可以抽象地写为 $\min_{\mathcal Q}D(\mathcal Q)$，约束是 $R(\mathcal Q)\leq R_0$，其中 $R$ 是平均存储比特数，$D$ 是明确指定的失真。这里可以用参数误差、输出误差或任务风险作失真，但不能在推导中随意替换它们。

固定长度编码中，每个索引使用 $b$ 位；若再进行熵编码，平均存储可能低于 $b$ 位，但随机访问和并行解码可能更困难。推理内核往往需要规则布局，所以文件压缩比不必等于运行时访存压缩比。有效位宽还应包括尺度、零点、码本、对齐填充和未量化参数，而不只计算主权重数组。这个视角将在第\ref{sec:quant-system}节用于解释为什么“四位”未必恰好是每参数四位。

\section{量化粒度、分组方式与混合精度}
\label{sec:quant-granularity}

\subsection{逐张量、逐通道和逐组}
\keyterm{逐张量量化（per-tensor quantization）}让整个张量共享一套尺度和零点，元数据少，算子组织简单。如果不同通道的范围差别很大，最大范围会决定全局间距，使小范围通道损失分辨率。\keyterm{逐通道量化（per-channel quantization）}为指定通道各自配置尺度。在线性层中，权重按输出通道量化意味着第 $i$ 行共享 $s_{w,i}$；卷积则通常可以把每个输出滤波器的权重看成一组。选择哪根轴，应该与数据分布和点积的计算结构共同决定。

\keyterm{逐组量化（group-wise quantization）}把权重沿约定方向分成固定大小的组。若每行沿输入维度每 $g$ 个权重共用一套参数，组索引为 $r$，就有 $\widehat W_{ij}=s_{i,r}(q_{ij}-z_{i,r})$。组大小 $g$ 越小，通常越能适配局部范围，但这只是经验趋势：尺度估计误差、优化策略、布局变化和不同端点约定都可能使单次结果不单调。若不允许组跨越某些矩阵边界，末尾不足一组时还会产生填充或特殊处理。

假设每组 $g$ 个权重，主编码每个 $b$ 位，尺度用 $b_s$ 位，零点用 $b_z$ 位，忽略填充时有效位宽为
\begin{equation}
 b_{\rm eff}=b+\frac{b_s+b_z}{g}.
 \label{eq:quant-effective-bits}
\end{equation}
例如四位权重、十六位尺度、十六位零点和组大小 $128$ 对应 $4.25$ 位每参数；若零点只占四位并能紧凑打包，则为 $4.15625$ 位。两种格式都可能叫“四位量化”，但存储成本不同。应使用实际序列化格式中的元数据位数，不能根据数学上零点只需少量编码就假定实现确实如此。

\subsection{激活的逐词元尺度和动态尺度}
语言模型的输入激活可整理为 $\mathbf X\in\mathbb R^{d\times N}$，每列对应一个词元位置。\keyterm{逐词元量化（per-token quantization）}为每列配置 $s_{x,n}$，使不同位置的整体幅度不互相挤占范围。它能够适配某个词元整体偏大，却不能自动解决同一个词元内部少数通道特别大的问题。逐通道激活量化则沿另一方向适配分布，可能更符合异常通道结构，却会给整数点积带来不同的尺度组织。

\keyterm{静态量化（static quantization）}通常在部署前用校准数据确定激活量化参数；\keyterm{动态量化（dynamic quantization）}在运行时根据当前输入计算某些尺度或范围。这一分类与 PTQ、QAT 是不同维度：PTQ 可以产生静态或动态激活方案，QAT 也可以模拟运行时动态量化。动态方案需要额外的归约、尺度计算和可能的同步，但在输入幅度变化较大时可能降低饱和风险。

动态量化不意味着每一个量化参数都在线更新。权重通常是固定的，可以离线打包；激活尺度可能按词元实时计算，而缓存尺度可能按新写入块计算后永久保存。应把具体张量的生命周期写清楚，避免用一个“动态”标签掩盖多种不同的执行过程。

\subsection{沿归约轴变化的尺度为何特殊}
若权重按输出行使用 $s_{w,i}$，激活按输入向量使用 $s_x$，则整个点积共用乘积尺度 $s_{w,i}s_x$，可以先累计整数乘积再一次重定标。若权重在同一行内分组，则
\begin{equation}
 \widehat y_i=\sum_r s_{w,i,r}s_x
 \sum_{j\in G_r}(q_{w,ij}-z_{w,i,r})(q_{x,j}-z_x)+\beta_i,
 \label{eq:quant-group-dot}
\end{equation}
其中 $G_r$ 是第 $r$ 组的输入坐标集合。不同组的部分和具有不同物理尺度，不能直接当成同一量纲的整数相加。实现需要分别重定标、转到共同累加尺度，或采用融合解码与乘加。

这说明最小误差的粒度不一定是最快粒度。量化设计应先考虑目标内核允许什么尺度广播方式、组大小和内存布局，再在兼容范围内比较误差。只有当两个候选方案都能落到真实内核上时，理论上的压缩与精度比较才具有直接部署意义。

\subsection{混合精度的资源分配}
\keyterm{混合精度量化（mixed-precision quantization）}为不同层、通道或元素分配不同精度。设第 $\ell$ 层选择位宽 $b_\ell$，成本为 $C_\ell(b_\ell)$，可把设计写成
\[
 \min_{b_1,\ldots,b_L}\mathcal E(b_1,\ldots,b_L),
 \qquad \sum_{\ell=1}^L C_\ell(b_\ell)\leq C_0.
\]
这里 $\mathcal E$ 应是实际关心的损失，$C_0$ 是预算。把 $\mathcal E$ 近似分解为各层敏感性的和，便于搜索，但忽略了层间误差耦合。只对单层做低精度实验，可以获得有用的诊断，却不能保证组合后的最优性。

少量异常通道保留较高精度是混合精度的一种形式。它可能显著降低损失，也可能引入额外索引、分支、矩阵拆分与内核启动。稀疏异常值的数量不是完整成本指标，还要看它们是否连续、是否容易共同访问，以及高低精度结果如何合并。

\section{从实数线性层到整数运算}
\label{sec:quant-integer}

\subsection{整数点积与零点修正}
先考虑一行权重与一个输入向量，量化重建为 $\widehat w_j=s_w(q_{w,j}-z_w)$、$\widehat x_j=s_x(q_{x,j}-z_x)$。代入内积得到
\begin{equation}
 \sum_{j=1}^d\widehat w_j\widehat x_j=s_ws_x A,
 \quad A=\sum_{j=1}^d(q_{w,j}-z_w)(q_{x,j}-z_x).
 \label{eq:quant-int-dot}
\end{equation}
对量化重建值而言，这是精确恒等式；与原始内积相比仍存在量化误差。展开累加器可得
\begin{equation}
 A=\sum_jq_{w,j}q_{x,j}-z_x\sum_jq_{w,j}
 -z_w\sum_jq_{x,j}+d z_wz_x.
 \label{eq:quant-zero-correction}
\end{equation}
因此整数矩阵乘法不是简单把浮点乘法中的数值强制转换为整数。零点项需要正确处理，某些权重求和可以离线预计算，而输入求和可能发生在运行时。若权重对称量化令 $z_w=0$，与输入求和有关的一项消失，便可简化运算。

对批量输入采用逐输出通道权重尺度和逐词元激活尺度时，$\widehat Y_{in}=s_{w,i}s_{x,n}A_{in}+\beta_i$。尺度矩阵是两个尺度向量的外积，这与整数矩阵乘法的输出结构容易对应。若尺度沿点积内部改变，就必须回到式\eqref{eq:quant-group-dot}的分组处理。

\subsection{偏置精度与重定标}
若采用整数偏置，通常令 $q_\beta=\operatorname{round}(\beta/(s_ws_x))$，从而把输出近似写成 $s_ws_x(A+q_\beta)$。在没有偏置饱和时，偏置引入的绝对误差至多为 $s_ws_x/2$。如果激活尺度 $s_x$ 在每次推理中改变，整数偏置的尺度也随之改变，此时不能无条件复用一个离线固定的 $q_\beta$；实现可以在较高精度中添加偏置，也可以运行时重定标。

设输出量化尺度为 $s_y$、零点为 $z_y$，则\keyterm{重定标（requantization）}为
\begin{equation}
 q_y=\operatorname{clip}\!\left(
 \operatorname{round}\!\left(M(A+q_\beta)\right)+z_y,
 q_{y,\min},q_{y,\max}\right),\qquad M=\frac{s_ws_x}{s_y}.
 \label{eq:quant-requant}
\end{equation}
这里的编码还会引入一次舍入和可能的饱和。为实现整数运算，可用整数乘子和移位近似 $M$，例如 $M\approx M_0/2^r$。需要根据范围选择中间乘积位宽，并实现约定的舍入规则；对于大于一的比例，还可能需要左移或不同的指数组织。不能默认所有重定标比例都小于一。

整数推理的经典构造可参考 Jacob 等人的工作\cite{jacob2018}。本文的关键结论是：量化、乘法、累加、偏置和输出重定标共同定义一个算子，训练时的模拟也必须与整条算术路径一致。只模拟权重舍入，却忽略输出饱和或偏置转换，可能高估部署精度。

\subsection{累加器为什么通常需要更宽}
若每个中心化整数满足 $|q_{w,j}-z_w|\leq R_w$、$|q_{x,j}-z_x|\leq R_x$，则 $|A|\leq dR_wR_x$。一个有符号 $p$ 位累加器的正端点是 $2^{p-1}-1$，因此充分条件为
\begin{equation}
 dR_wR_x+|q_\beta|\leq2^{p-1}-1.
 \label{eq:quant-overflow}
\end{equation}
若两个操作数都使用 $[-127,127]$ 的对称范围，三十二位有符号累加器在不含偏置的粗略最坏情形下要求 $d\cdot127^2\leq2^{31}-1$。非对称编码中心化后的范围可能达到 $255$，不能沿用 $127$ 的界。真实内核也可能分块累计或采用其他中间格式，需要逐层核对。

更多累加位不能恢复输入量化丢掉的信息，却可以避免误差再次被溢出放大。浮点累加同样存在舍入，而且不同归约顺序可能产生不同结果。这里的整数恒等式默认中间运算没有溢出；只有在这个条件下，整数加法顺序才不会改变数学结果。

\subsection{残差相加、激活函数和图边界}
设两个残差分支分别表示为 $a=s_a(q_a-z_a)$、$b=s_b(q_b-z_b)$。输出尺度为 $s_o$ 时，编码应按 $s_a/s_o$ 和 $s_b/s_o$ 分别转换后再合并，或在较高精度中相加后量化。只有在尺度和零点安排满足相应关系时，才可简化为直接的整数加法。忽略尺度差别相当于把不同计量单位的数值直接相加。

对仿射量化后的 ReLU，实数零对应编码 $z$，因此应执行 $q\mapsto\max(q,z)$，而不是总使用 $\max(q,0)$。GELU、softmax、归一化和除法等算子通常需要更复杂的近似或较高精度路径；是否全部量化应由误差收益与内核支持决定。计算图中频繁的量化与反量化转换还会消耗时间，所以融合边界与数值规则同样重要。

\section{校准：用有限数据确定量化参数}
\label{sec:quant-calibration}

\subsection{校准数据提供了什么信息}
\keyterm{校准（calibration）}是用代表性输入估计激活范围、通道统计量或输出敏感性，并选择量化参数的过程。它不一定需要人工标签，因为层重建可以把原模型输出作为目标。无标签不等于无信息：文本领域、语言、长度、提示格式和特殊词元都会改变激活分布。图像中的预处理、分辨率和归一化也会改变尺度，必须与部署路径相符。

最简单的观察器记录每个张量的最小值和最大值。它实现容易，却对单个极端值敏感，而且样本增加时观测范围往往扩大，未必使低位量化更精确。指数移动平均可以平滑统计变化，但它引入顺序依赖，也不能替代对长尾情况的覆盖。校准不是在一个小集合上寻求漂亮误差，而是估计未来执行中合适的数值策略。

权重量化若仅依据静态权重极值，可以不访问输入数据；若优化层输出误差，就需要输入的二阶信息或直接重建样本。因此“无数据量化”和“无标签量化”差别很大。某个方法可以只保留统计矩阵而丢弃原始样本，但统计量仍然来源于数据，不能据此称为没有校准信息。

\subsection{极值、百分位、均方误差和分布匹配}
极值校准直接覆盖观测范围；百分位校准选择指定分位点作为端点，容许少量饱和；均方误差校准搜索阈值，使 $N^{-1}\sum_i(Q_\alpha(x_i)-x_i)^2$ 尽量小。百分位方法控制的是样本比例，而不是饱和幅度，一个极远的异常值仍可能产生很大的平方损失。均方误差方法则可能为少量幅度很大的值牺牲大量中等数值。

还有方法对激活建立直方图，比较原分布与量化后重建分布，使用 KL 散度等准则选阈值。对离散概率 $p_k,r_k$，$D_{\rm KL}(p\|r)=\sum_{k:p_k>0}p_k\log(p_k/r_k)$，当 $p_k>0$ 而 $r_k=0$ 时为无穷大。因此实际算法需要明确如何合并桶、重分配质量和处理零概率。不同实现中的“熵校准”不必是同一个目标；它也不等于直接最小化模型输出概率的 KL 散度。

直方图桶数太少会掩盖分布细节，太多又容易产生大量空桶和不稳定估计。阈值搜索的计算量、候选范围与样本覆盖应一并记录。若最终只报告最优阈值，却不说明用什么数据和目标选择，就难以判断该阈值能否迁移到新的模型层或应用场景。

\subsection{逐层校准与真实输入分布}
设前面几层已经量化，当前层接收到的输入为 $\widehat{\mathbf X}$，而浮点网络产生的输入为 $\mathbf X$。只最小化 $\|\widehat{\mathbf W}\mathbf X-\mathbf W\mathbf X\|_F^2$，忽略了上游误差；比较 $\widehat{\mathbf W}\widehat{\mathbf X}$ 与浮点目标则能纳入已经累积的偏差。后者的目标仍可能受校准集限制，而且需要明确是在修正局部层误差还是整体前缀误差。

实践中的逐层流程通常先缓存有限输入或统计量，再量化一层，把量化后的输出送给后续层。这样内存需求可以受控，但顺序本身会影响结果。若把整个块作为重建单位，可以考虑残差和非线性之间的相互作用，代价是更复杂的优化和更多保存状态。重建粒度与量化粒度是不同概念：一个块内可以包含多种尺度和编码方式。

\subsection{校准集、验证集与测试集的分工}
校准集用于拟合范围或重建参数；验证集用于选择组大小、位宽和搜索超参数；测试集用于最终报告。不断根据测试集选择方案会使测试结果偏乐观。若数据有限，可以使用明确的数据切分或交叉验证，但应记录每种用途，不能把多次选优后的结果当作一次独立测试。

对语言模型，校准样本数与词元数不能混为一谈。一百个长度为几十的短句，与一百个长上下文，在激活覆盖、缓存规模和词元相关性上差别很大。应同时记录序列数、有效词元数、长度分布、语言和任务类型；填充词元是否参与统计也要明确。真实长上下文出现的异常位置，如果从未被采样，单纯扩大短文本数量可能无法补救。

\section{从局部误差到网络输出误差}
\label{sec:quant-propagation}

\subsection{线性层中的一阶项与交叉项}
令 $\widehat{\mathbf W}=\mathbf W+\mathbf E_W$、$\widehat{\mathbf x}=\mathbf x+\mathbf e_x$，偏置误差为 $\mathbf e_\beta$。则
\begin{equation}
 \widehat{\mathbf y}-\mathbf y
 =\mathbf E_W\mathbf x+\mathbf W\mathbf e_x
 +\mathbf E_W\mathbf e_x+\mathbf e_\beta.
 \label{eq:quant-linear-error}
\end{equation}
前三项依次是权重量化、激活量化及二者相互作用的贡献。只有当两类误差都足够小时，才有理由忽略交叉项。对应范数界为
\[
 \|\widehat{\mathbf y}-\mathbf y\|_2
 \leq\|\mathbf E_W\|_2\|\mathbf x\|_2
 +\|\mathbf W\|_2\|\mathbf e_x\|_2
 +\|\mathbf E_W\|_2\|\mathbf e_x\|_2+\|\mathbf e_\beta\|_2.
\]
这个界便于找出可能放大的因素，却未利用误差抵消和输入分布，因此通常不能作为实际精度下降的准确预测。

对于一行权重，如果输入坐标独立、均值为零，且把权重误差视为固定，则式\eqref{eq:quant-output-mse}简化为 $\sum_j\delta_j^2\operatorname{Var}(X_j)$。方差大的输入坐标对应更高的误差代价。若独立或零均值条件不成立，非对角项和均值项必须保留；不能因为公式更简单就默认删去。激活感知方法正是试图利用这种不均等的重要性。

\subsection{深层传播与 Lipschitz 上界}
假设第 $\ell$ 层原映射为 $f_\ell$，量化后映射为 $\widehat f_\ell$，实际输入和量化输入分别为 $\mathbf h_{\ell-1}$ 与 $\widehat{\mathbf h}_{\ell-1}$。若 $f_\ell$ 在相关区域具有 Lipschitz 常数 $L_\ell$，即 $\|f_\ell(u)-f_\ell(v)\|\leq L_\ell\|u-v\|$，并令
\[
 \eta_\ell=\|\widehat f_\ell(\widehat{\mathbf h}_{\ell-1})
 -f_\ell(\widehat{\mathbf h}_{\ell-1})\|,
\]
则误差递推为 $e_\ell\leq L_\ell e_{\ell-1}+\eta_\ell$，其中 $e_\ell=\|\widehat{\mathbf h}_\ell-\mathbf h_\ell\|$。若初始输入相同，反复展开可得
\begin{equation}
 e_L\leq\sum_{k=1}^L\eta_k\prod_{j=k+1}^L L_j.
 \label{eq:quant-deep-bound}
\end{equation}
这里空乘积取一。它解释了早期误差可能被后续层放大，但不意味着真实误差必然随深度指数增长；上界可能非常宽松。

对于残差映射 $f(h)=h+g(h)$，若 $g$ 的 Lipschitz 常数为 $L_g$，则可取上界 $1+L_g$。残差并不自动消除量化误差，只是改变了信息和误差的传播路径。归一化的敏感性则依赖局部方差与稳定常数，不能简单说“归一化总会压小误差”。在方差很小的区域，分母可能使某些扰动更敏感。

\subsection{分类边界与自回归分歧}
设原模型逻辑值最大类别为 $k_*$，它相对于其他类别的最小间隔为 $\gamma=a_{k_*}-\max_{j\ne k_*}a_j>0$。若量化后每个逻辑值误差的绝对值都不超过 $\varepsilon$，且 $2\varepsilon<\gamma$，则最大类别保持不变，因为任意竞争类别与原最大类别的差最多缩小 $2\varepsilon$。这是一个可直接验证的充分条件，不是必要条件，也没有保证概率完全不变。

对生成模型，每一步下一个词元分布的变化可能导致选择不同词元，之后两条生成路径面对的上下文已经不同。因此完整文本差异不能简单归结为某一层的均方误差。贪心解码对接近并列的逻辑值特别敏感，随机采样还受到随机数和温度影响。评估时应同时看固定上下文上的概率损失和自由生成任务，避免用单一文本是否逐字一致作为唯一质量标准。

\subsection{偏差修正的能力与局限}
若量化层的平均输出偏差是 $\mathbf m_e=\mathbb E[\widehat{\mathbf y}-\mathbf y]$，在偏置中减去其估计值可以校正均值。对于固定协方差下的平方误差，减去真实误差均值能去掉均方误差中的偏差平方项；它不能消除剩余方差和输入相关的误差。若校准分布与部署分布不同，修正值还可能不再合适。

偏差修正不应通过测试标签反复拟合。对于具有非线性和残差的网络，局部均值修正也可能改变后续激活分布，因此应重新验证完整模型。它是一种低成本补救手段，不能替代对错误量化轴、零点或累加溢出的修复。

\section{训练后量化的重建目标与二阶补偿}
\label{sec:quant-ptq}

\subsection{最近舍入为什么不是输出最优}
固定每个权重可用的码字集合，逐元素最近舍入最小化 $\|\widehat{\mathbf W}-\mathbf W\|_F^2$。但校准输出更自然的目标是
\begin{equation}
 \min_{\widehat{\mathbf W}\in\mathcal C^{m\times d}}
 \frac1N\|\widehat{\mathbf W}\mathbf X-\mathbf W\mathbf X\|_F^2.
 \label{eq:quant-layer-reconstruction}
\end{equation}
这里为简洁使用公共码本 $\mathcal C$，逐组码本时应把约束改为各组各自的集合。设误差矩阵为 $\boldsymbol\Delta$，目标等于 $\operatorname{tr}(\boldsymbol\Delta\widehat{\mathbf M}\boldsymbol\Delta^\top)$。只有当 $\widehat{\mathbf M}$ 与单位矩阵成比例等特殊条件成立时，这个目标才与逐元素平方误差具有相同的直接结构。

一个反例是单个输入 $\mathbf x=(1,1)^\top$，权重 $\mathbf w=(0.49,0.49)^\top$，码本取整数。最近舍入给出 $(0,0)$，输出误差为 $-0.98$；选 $(0,1)$ 时，权重平方误差略大，输出误差却只有 $0.02$。这不是说向上舍入总是更好，而是说明舍入决策之间存在耦合。用数据来选择舍入方向，可能比独立最小化参数误差更适合保持函数。

\subsection{自适应舍入与局部优化}
把权重写成 $w/s=\lfloor w/s\rfloor+r$ 后，可以为每个权重选择向下或向上舍入。所有选择组成离散优化问题，精确枚举的复杂度随权重数指数增长。一个可行思路是用连续变量 $h(v)\in[0,1]$ 暂时代替二元选择，优化重建误差并逐步鼓励 $h(v)$ 靠近零或一，最后恢复离散舍入。AdaRound 是这一路线的代表\cite{nagel2020}。

连续松弛让梯度优化成为可能，却没有把问题变成全局凸优化。初始化、正则权重、退火过程和重建数据都会影响结果。若只在很少的校准输入上优化，还可能过度适应这些输入的相关结构。实践中应保留简单最近舍入基线，并在独立输入上比较收益，才能知道复杂优化是否值得其时间成本。

\subsection{一个坐标固定时的最优连续补偿}
为了理解二阶补偿，考虑一行权重在当前可调整坐标上的扰动 $\boldsymbol\delta\in\mathbb R^r$。设局部代价为 $\tfrac12\boldsymbol\delta^\top\mathbf H\boldsymbol\delta$，其中 $\mathbf H$ 对称正定；现在必须把第 $j$ 个坐标改变量固定为 $\epsilon$，其他坐标仍可连续调整。问题是
\[
 \min_{\boldsymbol\delta}\tfrac12\boldsymbol\delta^\top\mathbf H\boldsymbol\delta,
 \qquad \mathbf e_j^\top\boldsymbol\delta=\epsilon,
\]
其中 $\mathbf e_j$ 为第 $j$ 个标准基向量。用拉格朗日乘子求解，驻点满足 $\mathbf H\boldsymbol\delta-\lambda\mathbf e_j=0$，代入约束得到
\begin{equation}
 \boldsymbol\delta^*=\frac{\epsilon}{(\mathbf H^{-1})_{jj}}\mathbf H^{-1}\mathbf e_j,
 \qquad
 \Delta L_{\rm quad}=\frac{\epsilon^2}{2(\mathbf H^{-1})_{jj}}.
 \label{eq:quant-compensation}
\end{equation}
正定性保证该驻点是约束下唯一最小点。这一结论对指定二次问题严格成立，展示了其他坐标如何分担固定坐标的误差；它没有保证后续所有坐标都被离散化后仍达到全局最优。

在线性重建中，可以取 $\mathbf H=2\mathbf X\mathbf X^\top/N$，因为平方目标的 Hessian 正是该矩阵。当样本不足或输入共线时，它可能奇异；采用 $\mathbf H_\lambda=2\mathbf X\mathbf X^\top/N+\lambda\mathbf I$、$\lambda>0$ 可以保证可逆，并相当于在目标中加入扰动范数惩罚。矩阵条件数过大时，数值求逆不稳定；实际实现宜用稳定分解和线性方程求解，而不是天真地重复显式求逆。

\subsection{GPTQ 的位置与逐步消元}
GPTQ 使用层重建和近似二阶信息，对逐步量化产生的误差进行补偿，并通过批量更新和矩阵分解改进大模型上的计算效率\cite{frantar2023}。可以把它理解为：先确定量化规则与输入统计，再固定一部分权重，把造成的误差转移到尚未固定的权重，最后完成整个离散化。这里的二阶信息主要对应局部重建目标，不能直接称为整个语言模型任务损失的精确 Hessian。

若当前 Hessian 的逆写为 $\mathbf A=\mathbf H^{-1}$，固定坐标 $j$ 后，剩余坐标集合 $F$ 对应的约束子问题需要使用
\begin{equation}
 (\mathbf H_{FF})^{-1}
 =\mathbf A_{FF}-\frac{\mathbf A_{Fj}\mathbf A_{jF}}{A_{jj}}.
 \label{eq:quant-inverse-update}
\end{equation}
这是分块矩阵求逆关系的结果。式\eqref{eq:quant-compensation}之后不能继续让已经固定的坐标自由变化，也不能一直复用未消元的完整逆矩阵，否则就不是同一个约束问题。分组尺度、量化顺序和数值阻尼都会影响具体算法，阅读实现时应区分核心数学与工程变体。

校准输入构造稠密 $d\times d$ 二阶矩约需 $O(Nd^2)$ 运算，保存矩阵需要 $O(d^2)$ 空间，朴素分解的代价约为 $O(d^3)$。这是相关核心步骤的量级说明，不是完整实现总复杂度或运行时间承诺。宽层、样本数、分块方法、内存转移和内核效率都会改变实际开销。局部补偿通常需要比最近舍入更多资源，所以应比较其精度收益与量化准备成本。

\section{量化感知训练与替代梯度}
\label{sec:quant-qat}

\subsection{伪量化的前向语义}
\keyterm{伪量化（fake quantization）}在浮点计算图中执行编码再解码，即把 $x$ 替换为 $Q(x)$，但常仍用浮点张量保存结果。它的作用是模拟量化后的数值，而不是直接实现存储压缩或整数加速。浮点矩阵乘法使用的输入虽然落在离散网格上，真实内核仍可能完全是浮点运算。因此，伪量化的速度不应作为最终量化模型的推理速度报告。

QAT 常保留一份较高精度的\keyterm{潜在权重（latent weight）} $\mathbf W$，前向用 $Q(\mathbf W)$，反向更新潜在权重。这样多个小于量化间距的梯度更新可以先在潜在权重中累积，直到跨越舍入边界才改变离散值。如果直接对已经离散的权重施加很小更新并立刻舍入，许多更新会永久消失。优化器状态通常也需要较高精度，这说明 QAT 的训练内存和推理内存必须分别核算。

\subsection{直通估计器为何是一种人为选择}
固定尺度和零点时，最近舍入函数在每个开区间内为常数，导数几乎处处为零，在跳变处不可导。直接使用其真实导数，梯度难以更新潜在权重。\keyterm{直通估计器（straight-through estimator, STE）}在反向中人为指定替代导数，例如对实际重建区间 $[\ell,u]$ 取
\begin{equation}
 \frac{\partial Q(x)}{\partial x}\;\mathrel{\mathop{\approx}^{\rm STE}}\;
 \mathbf 1_{\{\ell<x<u\}}.
 \label{eq:quant-ste}
\end{equation}
这里 $\mathbf 1$ 是指示函数；也有实现始终用一或采用不同裁剪规则。式中的近似符号表示算法定义的代理梯度，不是原离散函数导数的严格近似定理，也不自动无偏。

令 $\operatorname{sg}(v)$ 表示前向取值为 $v$、反向把它视作常量的停止梯度算子，表达式 $x+\operatorname{sg}(Q(x)-x)$ 在前向等于 $Q(x)$，对 $x$ 的反向导数则为一。若需要裁剪区间外导数为零，应进一步模拟裁剪的导数。不同伪量化模块可能在权重、尺度和零点上使用不同代理规则，复现时必须核对这些约定。

替代梯度可以在经验上找到更适合量化的参数，但它不保证每一步降低真实离散目标。学习率过大还可能导致潜在权重在阈值两侧振荡，离散权重不断翻转。检查训练过程时，应同时观察任务损失、饱和比例、尺度变化和码字切换，而不是只看潜在权重的梯度范数。

\subsection{可学习尺度的梯度推导}
考虑零点为零的简化伪量化 $\widehat v=s\,\operatorname{round}(\operatorname{clip}(v/s,q_{\min},q_{\max}))$，暂时固定整数端点。这里先裁剪归一化输入再舍入，与整数端点下先舍入再裁剪具有相同的前向数值；但定义反向规则时，需要明确选择对连续裁剪求导、对舍入使用 STE 导数一。在 $q_{\min}<v/s<q_{\max}$ 的内部区域，链式法则给出代理尺度梯度
\begin{equation}
 \frac{\partial\widehat v}{\partial s}\;\mathrel{\mathop{\approx}^{\rm STE}}\;
 \operatorname{round}(v/s)-v/s.
 \label{eq:quant-scale-gradient}
\end{equation}
若 $v/s<q_{\min}$，代理梯度为 $q_{\min}$；若 $v/s>q_{\max}$，代理梯度为 $q_{\max}$。边界处的约定需要具体指定，不能只根据舍入后的码字是否等于端点判断反向分支。这个推导说明，学习尺度不仅改变范围，也改变所有内部网格位置；它与单独训练一个裁剪阈值相关，却不能无条件视为相同参数化。

LSQ 方法把步长作为可学习参数，并对其梯度进行专门缩放\cite{esser2020}。实际优化中，尺度由很多元素共享，其梯度可能随元素数量显著变化，因此需要与权重学习率协调。为保证 $s>0$，可以用正值参数化、投影或下界限制。尺度接近零时的数值稳定性也需要处理，否则除法可能产生无穷大或非数值。

\subsection{一套完整的 QAT 训练流程}
给定浮点模型、训练数据、量化格式和部署算子约束，先确定需要插入伪量化的位置，并用短暂统计初始化范围。随后以前向伪量化、反向代理梯度训练潜在权重和允许学习的尺度；观察器若继续更新，其更新规则必须与尺度学习协调。训练后期通常需要固定某些统计量并降低学习率，让网络适应最终网格。何时冻结不是统一定理，应通过验证数据判断稳定性。

对于批量归一化，训练态使用当前小批量统计，推理态使用固定统计；两种前向不同。若最终部署要折叠归一化，应使训练末期的模拟与折叠后量化相匹配，否则训练中适应的范围可能与导出后不同。权重、尺度和归一化参数之间可能相互补偿，随意在导出时重新估计尺度也可能破坏训练取得的平衡。

最终导出要真正生成低位编码及元数据，并用部署内核执行代表性输入，比较其输出与冻结后的伪量化模型。即使二者使用同样位宽，也可能因偏置精度、舍入中点、零点修正和累加器溢出而不一致。QAT 只有在模拟与执行一致时，才算为真实目标算术做了训练。

\subsection{蒸馏、二值和三值网络}
可以在 QAT 中加入浮点教师与量化学生之间的输出匹配或中间特征匹配，以缓解低精度约束导致的学习困难。蒸馏目标提供额外监督，不会增加可表示码字，也不能保证弥补任意位宽降低。训练数据、教师质量和学生可表达能力仍然决定上限。若采用温度缩放的概率匹配，需要与任务损失权重共同调节，并在导出后评价实际网络。

\keyterm{二值量化（binary quantization）}常使用 $\{-a,a\}$。对固定向量 $\mathbf w\in\mathbb R^d$，若以平方误差拟合 $a\operatorname{sign}(\mathbf w)$，且 $a\geq0$，最优尺度为 $a=d^{-1}\sum_j|w_j|$；这是对 $\sum_j(|w_j|-a)^2$ 求导得到的结果。权重与激活都取 $\{-1,1\}$ 且使用适当位编码时，点积可写成 $2\operatorname{popcount}(\operatorname{XNOR}(\mathbf u,\mathbf v))-d$，其中 popcount 统计相同符号的个数。其效率还取决于位打包与支持的硬件指令。

\keyterm{三值量化（ternary quantization）}增加零码字，采用 $\{-a,0,a\}$。固定 $a>0$ 时，平方误差意义下，$|w|<a/2$ 的值更适合映射到零，超过阈值则映射到带符号的 $a$，中点按约定处理。若非零集合固定，最优 $a$ 是该集合内绝对值的平均。集合与尺度一起优化时仍需交替或搜索。三种状态的信息量下界与实际硬件存储不同，不能因为 $\log_2 3\approx1.585$ 就声称一个直接使用两位槽位的实现只占每参数 $1.585$ 位。

\section{等价变换：改变量化难度而保持浮点函数}
\label{sec:quant-transform}

\subsection{对角缩放的核心恒等式}
设 $\mathbf D=\operatorname{diag}(d_1,\ldots,d_d)$ 且所有 $d_j>0$。线性层满足
\begin{equation}
 \mathbf W\mathbf x=(\mathbf W\mathbf D)(\mathbf D^{-1}\mathbf x).
 \label{eq:quant-diagonal-transform}
\end{equation}
这是实数运算中的恒等式。若某个输入通道数值很大，可增大相应 $d_j$，使变换后的激活更小，同时让对应权重列更大。数值负担从激活一侧转移到权重一侧，而没有在浮点精确模型中改变函数。对两侧分别量化后，量化器一般不与变换交换，所以最终误差会随 $\mathbf D$ 改变。

例如标量乘积 $100\times0.01=1$，用 $d=10$ 可以改写成 $10\times0.1=1$。原来的两个操作数相差四个数量级，变换后相差两个数量级。这个例子只说明范围可重新分配；对于完整矩阵，各通道共享尺度且相互耦合，不能逐元素随意选择不同缩放而仍保留高效矩阵乘法。

\subsection{SmoothQuant 与 AWQ 的不同目标}
SmoothQuant 用通道缩放把部分激活异常值带来的量化困难转移到权重，面向权重和激活同时量化的场景\cite{xiao2023}。在本文矩阵方向下，若 $a_j$ 是校准激活第 $j$ 通道的最大绝对值，$w_j=\max_i|W_{ij}|$，一种代表性的平衡规则为
\begin{equation}
 d_j=\frac{a_j^\alpha}{w_j^{1-\alpha}},\qquad 0\leq\alpha\leq1,
 \label{eq:quant-smooth-scale}
\end{equation}
其中零值通道需要单独处理或加入正下界。增大 $\alpha$ 更倾向于缩小激活侧的幅度，但会增大权重侧负担；取 $\alpha=1/2$ 时，两侧相应最大幅度都成为 $\sqrt{a_jw_j}$。这是范围平衡的代数性质，不能推出任务损失在该点全局最小。

AWQ 则利用激活信息识别对权重量化更敏感的通道，通过等价缩放等步骤保护其有效分辨率，主要面向低位权重量化\cite{lin2024}。它的“激活感知”意味着用激活决定权重表示，不意味着推理时激活必然也是低位整数。若某个通道放大后被一整个组的尺度增长抵消，改善可能很有限，因此通道缩放必须结合分组与重建目标搜索，而不是简单认定放得越大越好。

两类方法都使用等价变换，但优化对象与目标执行格式不同。对相同模型比较时，应固定权重、激活和缓存的位宽、组大小及校准预算，才容易把方法差异与格式差异分开。不能把一个 W4A16 结果和一个 W8A8 结果的精度、速度差异全部归因于算法名称。

\subsection{异常值拆分与图上的变换约束}
另一条路线把输入通道拆分为普通集合 $R$ 和异常集合 $O$，从而精确地写成 $\mathbf W\mathbf x=\mathbf W_R\mathbf x_R+\mathbf W_O\mathbf x_O$。普通部分用低精度，异常部分用较高精度，再合并结果。LLM.int8() 是利用异常特征与混合精度矩阵乘法的代表\cite{dettmers2022}。其思想不是删除异常值，而是为它们分配不同算术路径；运行时异常值检测和拆分成本也应计入性能。

等价变换不能任意穿过非线性。对于正对角矩阵 $\mathbf D$，ReLU 满足 $\operatorname{ReLU}(\mathbf D\mathbf x)=\mathbf D\operatorname{ReLU}(\mathbf x)$，因为每个坐标只是乘正数；GELU 一般不满足同样的正齐次关系。带分支的计算图还要求所有消费者同步补偿，残差相加的两路也必须使用一致坐标。只修改一条路径，通常会在量化之前就改变原模型函数。

\subsection{正交旋转与异常值的分散}
若 $\mathbf R\in\mathbb R^{d\times d}$ 为正交矩阵，满足 $\mathbf R^\top\mathbf R=\mathbf I$，则
\begin{equation}
 \mathbf W\mathbf x=(\mathbf W\mathbf R^\top)(\mathbf R\mathbf x),
 \qquad\|\mathbf R\mathbf x\|_2=\|\mathbf x\|_2.
 \label{eq:quant-rotation}
\end{equation}
旋转保留欧氏长度，却会改变最大坐标和各通道分布。对 $\mathbf x=(M,0)^\top$，取 $\mathbf R=2^{-1/2}\begin{pmatrix}1&1\\1&-1\end{pmatrix}$，得到两个幅度为 $|M|/\sqrt2$ 的坐标。在共享最大值尺度时，这种分散可能改善普通坐标的表示。

对任意一个固定正交矩阵，也存在某些输入被旋转得更集中，因此不能宣称任意旋转都降低所有输入的量化误差。随机带符号的 Hadamard 变换常用于高效混合坐标；在适合的维度和实现下，快速变换可用 $O(d\log d)$ 次加减操作完成。QuaRot 将旋转思想用于大模型的多类张量量化\cite{ashkboos2024}，但具体可吸收哪些旋转，仍由归一化、激活、注意力和残差结构共同决定。

归一化对此尤其敏感。不含逐通道仿射系数的 RMS 归一化，因为分母只依赖平方和，可以与正交变换相容；带有逐通道系数时需额外吸收或补偿。LayerNorm 还减去坐标均值，一般正交矩阵不保持全一方向，所以不能直接套用同样结论。保证局部矩阵乘法等价，不等于已经保证整个 Transformer 等价。

\subsection{批量归一化折叠与校准顺序}
对输出 $u_i=\mathbf w_i^\top\mathbf x+\beta_i$，推理态批量归一化使用固定均值 $\mu_i$、方差 $\sigma_i^2$、缩放 $\gamma_i$、平移 $\eta_i$ 和稳定常数 $\varepsilon>0$。它可写成新的线性层
\[
 \mathbf w_i'=\frac{\gamma_i}{\sqrt{\sigma_i^2+\varepsilon}}\mathbf w_i,
 \qquad
 \beta_i'=\frac{\gamma_i(\beta_i-\mu_i)}{\sqrt{\sigma_i^2+\varepsilon}}+\eta_i.
\]
这是固定统计下的代数折叠，不适用于仍随小批量变化的训练态统计。折叠可能显著改变不同输出通道的权重范围，因此通常应针对实际部署图校准，而不是在旧权重上校准后再随意折叠。逐通道量化往往有助于适配这种通道范围差异。

\section{量化基础模型上的低秩适配与非均匀四位表示}
\label{sec:quant-qlora}

\subsection{冻结量化权重与可训练适配器}
设冻结的重建基础权重为 $\widehat{\mathbf W}_0\in\mathbb R^{m\times d}$，低秩适配器为 $\mathbf B\mathbf A$，其中 $\mathbf B\in\mathbb R^{m\times r}$、$\mathbf A\in\mathbb R^{r\times d}$ 且 $r\ll\min(m,d)$。忽略偏置时，前向为
\begin{equation}
 \mathbf y=\widehat{\mathbf W}_0\mathbf x+\rho\mathbf B\mathbf A\mathbf x,
 \label{eq:quant-lora}
\end{equation}
其中 $\rho$ 是适配器缩放系数。训练只更新 $\mathbf A,\mathbf B$，基础权重编码保持冻结。QLoRA 将量化基础模型与低秩适配结合，并使用 NF4、量化尺度等手段降低微调内存\cite{dettmers2023}。

反向传播仍需要通过基础线性映射把梯度传给更早的激活，因为 $\partial\mathbf y/\partial\mathbf x$ 包含 $\widehat{\mathbf W}_0$。但不必计算或保存基础权重的参数梯度与优化器状态。这个区别解释了“梯度经过量化模型”与“量化基础权重本身被训练”为什么不是一回事。适配器、激活和矩阵乘法通常还使用更高精度，所以四位基础权重也不表示所有训练计算都在四位上进行。

\subsection{NF4 是码本而非四位整数网格}
NF4 是面向近似正态权重设计的非均匀四位码本，编码保存十六个重建值之一的索引，通常配合分块归一化。它在零附近分配较密的重建值，在尾部较稀；其具体码字构造应以原论文和实现为准。NF4 不等于通用 IEEE 浮点格式，也不等于等间距 INT4。码字依赖特定分布设计，并不能推出它对所有权重、激活或所有损失都最优。

要特别区分分位点设计与第\ref{sec:quant-distortion}节的平方误差最优质心条件。按概率分布安排码字是一种有动机的设计，平方误差最优码字则还应与其最近邻区域质心相匹配，二者不由名称自动等价。此外，把一个有限块除以该块最大绝对值，会改变标准化后元素的联合分布；即使原始元素近似正态，归一化后的值也不是未经改变的独立标准正态样本。

\subsection{尺度再量化及其误差}
\keyterm{双重量化（double quantization）}在这里指进一步量化第一层量化所用的尺度等常数，并非把同一张量无目的地量化两遍。设一个组的重建权重为 $\widehat w=s c_k$，尺度再量化为 $\widehat s=s+\delta s$，则新的重建值变成 $\widetilde w=\widehat s c_k$，额外误差精确为 $\delta s\,c_k$。同组权重通过共同的尺度误差 $\delta s$ 耦合；即使另加随机模型分析，也不能未经论证就把这些误差视为相互独立的白噪声。

若每 $g$ 个权重原本保存一个三十二位尺度，其开销为 $32/g$ 位每参数。若这些尺度再被编码为八位，并每 $h$ 个一级尺度共用一个三十二位二级尺度，忽略其他元数据时开销成为 $8/g+32/(gh)$。这是特定假设下的成本公式；是否还需要偏移量、对齐或更高精度异常尺度，应依据实际格式补充。

\subsection{合并、重新量化与部署误差}
把适配器合并到基础权重，数学上得到 $\widehat{\mathbf W}_0+\rho\mathbf B\mathbf A$。如果再使用低位格式保存，得到的是 $Q(\widehat{\mathbf W}_0+\rho\mathbf B\mathbf A)$，通常不等于原来两条计算路径的和。量化是非线性操作，不能把合并后的再次舍入当作无损导出。合并还可能改变每组最大值与尺度，从而影响组内大量权重。

若保留独立适配器，额外矩阵乘法和适配器存储需要计入推理成本；若合并再量化，则需要重新验证质量。低秩适配可以在任务数据上修正部分量化影响，但低秩空间未必覆盖全部量化误差方向。它也可能主要学习新任务信息，而非专门恢复原模型函数，评价时应说明训练目标与比较基线。

\section{Transformer 与键值缓存量化}
\label{sec:quant-kv}

\subsection{注意力计算中的两类误差}
对一个注意力头，设历史长度为 $T$，键向量组成 $\mathbf K\in\mathbb R^{T\times d_k}$，值向量组成 $\mathbf V\in\mathbb R^{T\times d_v}$，当前查询为 $\mathbf q\in\mathbb R^{d_k}$。省略已处理的掩码位置后，计算为
\begin{equation}
 \mathbf a=\frac{\mathbf K\mathbf q}{\sqrt{d_k}},\qquad
 \mathbf p=\operatorname{softmax}(\mathbf a),\qquad
 \mathbf o=\mathbf V^\top\mathbf p.
 \label{eq:quant-attention}
\end{equation}
键量化先影响分数 $\mathbf a$，再通过 softmax 改变对所有位置的权重；值量化在权重给定时直接改变加权和。因此键和值即使维度类似，也不承担相同数值角色，不能仅凭一张张量直方图就假定它们应使用完全相同的量化粒度。

若键误差为 $\mathbf E_K$，查询误差为 $\mathbf e_q$，分数误差为 $(\mathbf E_K\mathbf q+\mathbf K\mathbf e_q+\mathbf E_K\mathbf e_q)/\sqrt{d_k}$。对于较小分数扰动，softmax Jacobian 为 $\mathbf J=\operatorname{diag}(\mathbf p)-\mathbf p\mathbf p^\top$，故 $\Delta\mathbf p\approx\mathbf J\Delta\mathbf a$。这是局部一阶近似；当关键分数次序被改变或扰动较大时，需要直接计算新分布。

若值误差为 $\mathbf E_V$，输出差精确分解为 $\mathbf V^\top\Delta\mathbf p+\mathbf E_V^\top\mathbf p+\mathbf E_V^\top\Delta\mathbf p$。固定注意力权重时，值误差是一个加权平均；但当权重也发生变化，交叉项不可无条件忽略。注意力中的异常值既可能来自输入表示，也可能来自位置变换与特定上下文，校准时应覆盖实际缓存写入位置。

\subsection{缓存容量与长上下文}
\keyterm{键值缓存（key--value cache, KV cache）}保存历史词元的键和值，避免每一步生成都重新计算它们。假设批量为 $B$，层数为 $L$，历史长度为 $T$，键值头数为 $H_{\rm kv}$，每头键和值维度都为 $d_h$，每元素保存 $b_{\rm kv}$ 位，则理想有效载荷为
\begin{equation}
 M_{\rm KV}=\frac{2BLTH_{\rm kv}d_hb_{\rm kv}}8\quad\text{字节}.
 \label{eq:quant-kv-memory}
\end{equation}
因子二对应键和值。分组查询注意力中，应使用实际键值头数而非查询头数；共享缓存、滑动窗口和不同层结构则需要相应修改公式。尺度、零点、页管理和未量化尾部另计。

例如取 $B=1$、$L=32$、$T=8192$、$H_{\rm kv}=8$、$d_h=128$，十六位缓存的理想载荷是 $1\,073\,741\,824$ 字节，即 $1$ GiB；四位缓存的理想载荷是 $0.25$ GiB。这里 GiB 等于 $2^{30}$ 字节。权重固定时，缓存仍随上下文和并发线性增加，因此只量化权重不一定解决长上下文服务的主要内存瓶颈。

\subsection{在线写入、分组和高精度尾部}
权重可以离线一次量化，缓存却随着生成不断追加。若每组跨越若干词元，组未写满时尚不知道未来最大值；不断改变尺度会要求重新量化已经保存的数据。可行设计包括固定离线范围、每个词元独立尺度、按块完成后再量化，或保留最近一小段较高精度缓存。每种选择都在误差、元数据、写入成本与读取效率之间取舍。

KIVI 根据键和值的分布特点使用不同量化轴，并结合残留的较高精度部分组织缓存\cite{liu2024}。这一设计说明缓存量化需要考虑数据布局和在线生命周期，而不是把权重量化器直接搬过去。不同模型、位置编码和注意力头结构可能改变分布，具体粒度是否适合仍应由实际统计与任务验证决定。

如果旧缓存每次读取后都反量化成完整高精度副本并长期保存，实际内存收益会被抵消。融合内核可以按需解码小块并参与注意力计算，但其效率取决于布局和算子支持。所谓缓存四位化应明确量化发生在存储、读取和计算的哪一环，尤其要检查峰值临时内存。

\subsection{长序列与生成行为的评估}
缓存量化的误差会影响后续每一步查询对历史的读取；关键事实位于很早位置时，效果可能与短文本困惑度不同。应在多种长度上检查检索、复制、对话历史利用和长程依赖任务，并区分位置效应与总长度效应。单次生成失败可以用于诊断，却不能替代足够数量样本的统计结果。

旋转位置编码等操作会改变键的坐标分布，因此必须记录缓存是在位置变换前还是之后被量化。对同一实数张量施加可逆变换再量化，一般与先量化再变换不同。与第\ref{sec:quant-transform}节相同，量化位置属于算法定义的一部分，不是可以随意调整的实现细节。

\section{存储、带宽与真实推理性能}
\label{sec:quant-system}

\subsection{模型文件、常驻内存和峰值内存}
若模型包含 $P$ 个量化参数，主编码为 $b$ 位，则理想载荷为 $Pb/8$ 字节。假设某模型有七十亿个参数，十六位权重载荷为 $14$ GB，四位载荷为 $3.5$ GB，这里的 GB 按 $10^9$ 字节计算；换成 GiB 数值会不同。实际模型还要保存量化元数据、未量化的嵌入或归一化参数，以及可能的张量对齐填充，因此文件大小不能直接用理想载荷替代。

推理常驻内存还包括缓存、运行时工作区、分配器保留空间和通信缓冲。加载时可能先读取浮点权重再转换，从而产生高于稳态的峰值；合并适配器或重排权重也可能临时保留多份副本。报告内存时应区分模型文件大小、加载峰值、推理峰值和有效缓存容量。同一个压缩格式在不同运行时中可能有很不相同的峰值表现。

对于混合精度模型，可以逐张量累加 $P_\ell b_\ell/8$，再加上各类元数据。若有部分层完全保留十六位，应明确其参数占比。尤其在小模型或很小的分组中，尺度和码本开销可能不可忽略；在大批量长上下文中，权重可能反而不是内存最大的组成部分。

\subsection{计算强度与带宽瓶颈}
设某次算子执行需要 $F$ 次运算、传输 $D$ 字节，\keyterm{算术强度（arithmetic intensity）}为 $I=F/D$。若设备相应内核的可达计算速率为 $P_{\rm comp}$，可达带宽为 $B_{\rm mem}$，理想重叠下的粗略时间下界为
\begin{equation}
 T\gtrsim\max\left(\frac{F}{P_{\rm comp}},\frac{D}{B_{\rm mem}}\right).
 \label{eq:quant-roofline}
\end{equation}
这是帮助定位瓶颈的简化模型，忽略启动、同步、缓存命中和调度等细节。若主要受带宽限制，减少权重读取字节可能带来收益；若主要受计算限制，则还需要低精度算术内核或更好的复用，单纯文件变小未必足够。

语言模型的\keyterm{预填充（prefill）}同时处理提示中的多个词元，往往能复用权重并形成较大的矩阵乘法；\keyterm{解码（decode）}逐步生成新词元，小批量时权重读取占比可能更高。这个趋势依赖具体模型、批量和硬件，并非固定定律。四位权重可能改善某一解码负载，却在另一种大批量预填充负载中受解码开销或计算内核限制，二者不矛盾。

\subsection{为什么四倍压缩不是四倍加速}
量化后的总时间可概念性地拆成数据传输、解包、尺度计算、乘加、重定标和调度等部分，其中部分过程可以融合或重叠。主权重字节数下降四倍，不会使所有部分同比下降。若只有原耗时比例 $f$ 的部分得到加速，且该部分加速倍数为 $r$，理想整体加速比是
\begin{equation}
 S=\frac1{(1-f)+f/r}.
 \label{eq:quant-amdahl}
\end{equation}
例如 $f=0.8$、$r=4$ 时，整体加速仅为 $2.5$ 倍；若还增加转换开销，实际值可能更低。这是时间构成上的推算，不是某个量化系统的实测结果。

四位权重常把两个码字放进一个字节，读取时需要按位提取、符号解释或零点调整。若每个元素都单独查表并产生不规则访问，解码代价可能抵消带宽收益。高效内核往往把解码、尺度乘法和矩阵乘加融合在寄存器或片上存储中。因而一个格式的优劣还取决于其打包布局是否适合目标硬件的向量化与矩阵指令。

\subsection{导出格式与执行契约}
导出文件应包含足够信息，使另一个程序明确恢复量化算子：整数是否有符号、码字排列、组大小、量化轴、尺度与零点的类型、末组填充方式、舍入规则，以及哪些层保留高精度。名义上同为 W4A16 的两个文件若这些约定不同，不能直接互换。模型配置中的量化名称只是入口，具体数值契约才决定正确性。

若后端不支持某一算子，运行时可能回退到高精度路径，并在图中插入额外转换。模型可以正确运行，却没有预期加速。检查性能时应记录实际内核和回退比例；比较精度时应确认执行的确是预定低精度模型，而不是某些敏感层被自动恢复为高精度。正确性与性能都要围绕真实执行图验证。

\section{贯穿编码、累加和输出的完整算例}
\label{sec:quant-example}

\subsection{原始线性函数与指定量化参数}
考虑一个两输入单输出线性层，权重 $\mathbf w=(0.34,-0.57)^\top$，输入 $\mathbf x=(0.8,1.3)^\top$，偏置 $\beta=0.10$。浮点实数模型的输出为
\[
 y=0.34\times0.8-0.57\times1.3+0.10=-0.369.
\]
为便于手算，指定权重尺度 $s_w=0.1$、零点 $z_w=0$、编码范围 $[-7,7]$；指定激活尺度 $s_x=0.1$、零点 $z_x=3$、编码范围 $[0,15]$。激活的实际重建范围是 $[-0.3,1.2]$，所以第二个输入将发生饱和。这里的尺度是为展示机制人为指定，不声称来自最优校准。

权重编码是 $(3,-6)$，重建为 $(0.3,-0.6)$；激活编码先计算 $(11,16)$，再裁剪为 $(11,15)$，重建为 $(0.8,1.2)$。权重误差为 $\mathbf e_w=(-0.04,-0.03)^\top$，激活误差为 $\mathbf e_x=(0,-0.1)^\top$。低位网格一方面引入舍入误差，另一方面由于预设范围不覆盖输入而产生饱和误差。

\subsection{整数累加、偏置与输出编码}
中心化激活编码为 $(8,12)$，因此整数点积为 $A=3\times8+(-6)\times12=-48$。偏置尺度为 $s_ws_x=0.01$，恰好有 $q_\beta=10$，无需额外偏置舍入。累加后得到 $A+q_\beta=-38$，重建输出为 $\widehat y=-0.38$。若忘记激活零点，直接计算 $3\times11-6\times15=-57$，就会得到错误的数值，这说明零点修正不可省略。

再指定输出尺度 $s_y=0.02$、零点 $z_y=32$、无符号八位编码范围 $[0,255]$，则重定标比例为 $M=0.5$，输出编码为 $q_y=-19+32=13$，解码仍为 $0.02(13-32)=-0.38$。本例输出恰好落在网格上，因此最终重定标未增加误差；一般情况下不会如此幸运。

将误差按式\eqref{eq:quant-linear-error}分解，有
\[
 \mathbf e_w^\top\mathbf x=-0.071,\qquad
 \mathbf w^\top\mathbf e_x=0.057,\qquad
 \mathbf e_w^\top\mathbf e_x=0.003.
\]
三项相加为 $-0.011$，与 $-0.38-(-0.369)$ 完全一致。虽然激活饱和误差不小，它与权重误差在这个输入上部分抵消；这种偶然抵消不应被当作校准范围合理的证据，应在多个输入上评价。

\subsection{二阶补偿的二维手算}
令局部二次代价为 $\tfrac12\boldsymbol\delta^\top\mathbf H\boldsymbol\delta$，其中 $\mathbf H=\begin{pmatrix}2&1\\1&2\end{pmatrix}$。现在把第一个坐标固定为 $\delta_1=0.1$。由 $\mathbf H^{-1}=\tfrac13\begin{pmatrix}2&-1\\-1&2\end{pmatrix}$ 和式\eqref{eq:quant-compensation}，最优连续补偿为 $\boldsymbol\delta=(0.1,-0.05)^\top$，代价为 $0.0075$。若不补偿第二坐标，取 $(0.1,0)^\top$，代价是 $0.01$。

若第二坐标最终也只能取某些离散增量，$-0.05$ 未必可实现。将它再量化后可能重新增加损失，所以这里证明的是一次固定坐标后的连续条件最优，并非完整离散问题最优。这个小例子把二阶信息的作用具体化：非对角项记录误差方向之间的关系，其他坐标可以沿代价较低的方向调整。

\subsection{不同误差指标可能给出不同排序}
设两种方法对同一组样本分别产生误差 $(0,0,0,1)$ 与 $(0.4,0.4,0.4,0.4)$。前者平均绝对误差是 $0.25$，后者为 $0.4$；前者均方误差是 $0.25$，后者为 $0.16$。平均绝对误差认为前者更好，均方误差认为后者更好。若第一个方案的大误差发生在关键决策边界附近，其任务风险还可能更高。

因此评估量化时不能只挑选最有利的一个指标。元素均方误差、层输出相对误差、最终交叉熵和任务成功率应按用途解释；它们之间存在联系，但不构成普遍单调关系。完整算例的价值在于检查数值路径，而最终质量仍要通过代表性任务验证。

\section{实现、诊断与实验设计}
\label{sec:quant-practice}

\subsection{从参考量化器到完整模型}
一个最小标量参考实现应把编码与解码分开，并显式保存参数。以下伪代码中的 round 必须采用指定中点规则，clip 必须是饱和操作，所有数学输入均假定有限；含非数值输入应在上游诊断，而不是默默依靠类型转换处理。
\begin{verbatim}
encode(x, scale, zero, qmin, qmax):
    assert scale > 0
    q = round(x / scale) + zero
    return clip(q, qmin, qmax)

decode(q, scale, zero):
    return scale * (q - zero)
\end{verbatim}
对张量实现，还需明确尺度广播轴；对真实低位格式，还需添加位打包与解包。这段参考代码用于验证数值定义，没有表达优化内核，也不意味着返回的整数数组已经按目标位宽紧凑保存。

一套基本 PTQ 流程可以按如下顺序组织，步骤中的比较都使用已预先划定的数据。
\begin{enumerate}
 \item 固定浮点模型、预处理、推理模式和评价协议，记录浮点输出与资源基线。
 \item 根据目标后端确定支持的格式、粒度、累加器和算子边界，并完成需要的推理图折叠。
 \item 运行校准输入，收集范围、二阶矩或层重建样本，明确有效词元和掩码处理。
 \item 先生成最近舍入基线，再按需要加入裁剪、分组、缩放或二阶补偿，使用验证集选择超参数。
 \item 用伪量化或参考解码图检查误差来源，定位敏感层与饱和通道，并只对有依据的问题调整方案。
 \item 导出低位编码和完整元数据，比较部署内核与参考图的数值结果，排查舍入和布局不一致。
 \item 在固定硬件及负载上测量质量、延迟、吞吐量和峰值内存，最后用保留测试集报告结果。
\end{enumerate}
流程结束条件应具体，例如达到预先约定的质量损失上限与内存预算，或在固定搜索预算内找到最优候选。若没有达到目标，保留较高精度或改变格式是有效结论，不应为了满足名义位宽而隐藏严重质量下降。

\subsection{数值正确性的分层检查}
标量层面检查零是否精确表示、网格点是否保持不变、越界是否饱和、中点规则是否一致，以及全零范围是否生成正尺度。张量层面检查逐通道与逐组的广播方向、末组处理、打包后解包是否恢复相同编码。算子层面检查整数点积与解码后浮点点积是否在预定误差内一致，偏置尺度和残差尺度是否正确。

对于整数累计，在没有溢出且未引入近似重定标的阶段，可以要求精确相等；对于浮点解码与不同归约顺序，应依据累加格式设置合理容差。容差过宽会掩盖错误，过窄又可能把正常浮点差异误判为实现失败。测试应覆盖负数、非零零点、端点、较长归约维度和不同形状，而不只是随机生成幅度很小的输入。

若某层误差突然远大于相邻层，首先检查轴、转置、零点和数据类型，而不是立即更换复杂算法。若所有层都有小幅相同方向偏差，可以检查舍入规则、偏置和范围选择。若参考图正常而真实内核失败，优先定位打包、布局和不支持算子；若短上下文正常而长上下文恶化，则进一步检查缓存和校准覆盖。诊断应让观测现象对应具体假设。

\subsection{质量指标与语言模型困惑度}
层级指标可以记录 $\|\widehat{\mathbf Y}-\mathbf Y\|_F^2$、相对误差 $\|\widehat{\mathbf Y}-\mathbf Y\|_F/(\|\mathbf Y\|_F+\epsilon)$、最大绝对误差和饱和比例，其中 $\epsilon>0$ 防止分母为零。余弦相似度可以描述方向，但不反映整体缩放；均方误差受输出尺度影响，跨层比较时应说明是否归一化。任何单一层级指标都只是任务质量的代理。

语言模型常使用\keyterm{困惑度（perplexity）}，对规定协议下的 $N$ 个目标词元定义为
\begin{equation}
 \operatorname{PPL}=\exp\!\left(-\frac1N\sum_{t=1}^N
 \log p(x_t\mid x_{<t})\right).
 \label{eq:quant-ppl}
\end{equation}
词元化器、上下文窗口、样本拼接、首词元处理和掩码都会影响数值。只有协议一致时，浮点与量化结果的差值才容易解释。困惑度相近并不保证代码执行、数学推理、长文本检索等任务成功率相近，应根据应用增加直接任务指标。

若比较多个随机种子或数据抽样，应报告均值及相应波动，而不是只报最好一次。对同一测试样本成对比较浮点与量化结果，通常更便于定位退化发生在哪类输入。小幅差异是否稳定，需要结合样本量和测量方差判断；一次小数点后的变化不必然具有实际意义。

\subsection{速度和内存测量的公平性}
推理计时前需要预热，并使编译、权重加载和内核选择等一次性成本与稳定执行成本分开。GPU 运算通常异步提交，计时要使用合适事件或同步边界，否则测到的可能只是提交任务的时间。批量、输入长度、输出长度、并发、采样策略、缓存状态和设备功率条件都应尽量固定。

\keyterm{首词元延迟（time to first token）}主要反映接收请求到产生首个输出的等待与预填充；\keyterm{每输出词元时间（time per output token）}反映后续解码速度；吞吐量则取决于并发与调度。一个方案可能提高总吞吐量却增加单请求等待，因此应按照目标服务指标判断。报告峰值内存时也要说明测量范围是否包含加载、预热和最大长度缓存。

格式对比至少应固定模型权重来源、任务协议和硬件，再列出位宽、量化轴、组大小、校准预算、是否使用适配器、内核版本及高精度回退。若一种方法使用更多数据或更长重建时间，应将这种成本纳入比较。本文只给出实验设计，没有在真实模型或设备上执行这些实验，也不提供未经测量的加速或精度数字。

\subsection{选择方法时的问题顺序}
先确定瓶颈：模型是否装不下，解码是否受权重带宽限制，长上下文是否被缓存占满，还是预填充计算量过大。再检查后端支持：需要什么低位内核、允许怎样的分组和尺度。随后选取简单基线，依据误差诊断增加校准、缩放或重建；只有 PTQ 无法满足质量要求且训练资源允许时，再考虑 QAT 或适配器微调。

问题顺序之所以重要，是因为优化不相关的对象可能毫无收益。缓存主导的场景只压缩权重，可能释放不了所需内存；计算主导的场景采用仅减少读取字节而增加解码运算的格式，可能更慢。研究方法提供的是工具，最终选择应由目标约束和可测量结果决定。

\section{练习、解答与进一步阅读}
\label{sec:quant-exercises}

\subsection{练习一：范围、零点和舍入}
使用无符号八位编码表示区间 $[-2,3.1]$，按式\eqref{eq:quant-minmax}计算尺度和零点，并求 $x=0.37$ 与 $x=3.4$ 的编码、重建值和误差。进一步说明，若把整数编码 $q=100$ 从这一张量复制到尺度为 $0.1$、零点为零的另一个张量中，实数含义是否保持不变。

\begin{solve}
区间长度为 $5.1$，故 $s=5.1/255=0.02$，$z=100$。$0.37/0.02=18.5$ 恰在中点，若采用最近取偶数，则编码为 $118$，重建 $0.36$，误差 $-0.01$；若采用中点远离零，编码为 $119$，重建 $0.38$，误差 $0.01$。$3.4$ 对应未裁剪编码 $270$，饱和后为 $255$，重建 $3.1$，误差 $-0.3$。原张量中的编码 $100$ 表示零，目标张量中的编码 $100$ 表示十，所以直接复制编码不保持数值。这里中点分析使用精确实数，浮点实现还须考虑十进制数的二进制近似。
\end{solve}

\subsection{练习二：元数据和缓存容量}
一个含十亿个量化权重的数组采用四位编码，每组六十四个权重，共享一个十六位尺度和一个八位零点，忽略填充与其他层。计算有效位宽及字节数。若式\eqref{eq:quant-kv-memory}中的批量由一变为四，同时上下文长度减半，缓存有效载荷变成原来的多少倍？

\begin{solve}
每参数位数是 $4+(16+8)/64=4.375$，总字节数为 $10^9\times4.375/8=546\,875\,000$，即约 $0.547$ GB。直接按四位计算只得到 $500\,000\,000$ 字节，会漏掉元数据。其他缓存参数不变时，批量扩大四倍、长度减半，总载荷成为原来的两倍；量化尺度等额外开销若不严格线性，还需另行计算。
\end{solve}

\subsection{练习三：相同参数误差为何不同}
设输入二阶矩为 $\mathbf M=\operatorname{diag}(100,1)$。两个候选权重误差分别为 $\boldsymbol\delta_a=(0.1,0)^\top$ 和 $\boldsymbol\delta_b=(0,0.1)^\top$。比较参数平方误差与输出均方误差，并说明这对按通道分配精度有什么启示。

\begin{solve}
两者参数平方误差都是 $0.01$，但输出均方误差分别为 $1$ 与 $0.01$。第一个输入方向的二阶能量更大，所以同样大小的权重误差在该方向上代价更高。可优先改善对应通道的量化精度，但这里的结论依赖给定输入统计；若部署分布改变或其他层补偿显著，仍需重新验证。
\end{solve}

\subsection{练习四：随机舍入与非线性}
把 $x=0.3$ 随机舍入到 $\{0,1\}$，使舍入无偏。求舍入值的方差，并比较先平方再求期望与原数的平方。这个例子是否支持“无偏量化不改变网络输出均值”？

\begin{solve}
应以概率 $0.3$ 输出一，以概率 $0.7$ 输出零，因此均值为 $0.3$，方差为 $0.3\times0.7=0.21$。因为输出只能取零或一，其平方期望仍为 $0.3$，而原数平方为 $0.09$。差额恰好等于方差。这反驳了上述无条件说法：无偏性可通过线性映射保留，但通常不能通过非线性映射保留。
\end{solve}

\subsection{练习五：缩放与旋转的条件}
证明正对角矩阵与 ReLU 可交换，并举例说明一般旋转与 ReLU 不可交换。再解释为什么只有线性层等价，还不足以证明整个模型在插入旋转后等价。

\begin{solve}
逐坐标有 $\max(d_jx_j,0)=d_j\max(x_j,0)$，因为 $d_j>0$。取 $\mathbf x=(1,-1)^\top$ 和第\ref{sec:quant-transform}节的二维矩阵 $\mathbf R$，则 $\mathbf R\mathbf x=(0,\sqrt2)^\top$，其 ReLU 不变；而 $\mathbf R\operatorname{ReLU}(\mathbf x)=(1/\sqrt2,1/\sqrt2)^\top$，二者不同。整个模型还包含非线性、归一化、位置变换和分支，必须为每个相关算子建立合法的补偿或不变性，局部乘法恒等式不足以覆盖这些操作。
\end{solve}

\subsection{练习六：梯度、冻结与再次量化}
在式\eqref{eq:quant-lora}中冻结基础权重，只训练低秩适配器。解释是否仍需向输入传播梯度，以及为什么 $Q(\widehat{\mathbf W}_0+\rho\mathbf B\mathbf A)$ 与原计算不一定相同。若训练后模型质量提高，能否据此断言量化误差全部被低秩矩阵消除？

\begin{solve}
若更早的层需要训练，就仍需通过 $\widehat{\mathbf W}_0$ 及适配器计算输入梯度；冻结只取消基础权重自身的参数更新。合并后再次量化会重新舍入并可能改变组尺度，因此通常产生新的误差。任务质量提高只能说明训练后的整体函数在该评价上更好，既可能来自补偿，也可能来自新任务学习；低秩矩阵的表达空间有限，不能由一个质量指标推出所有参数或函数误差都被消除。
\end{solve}

\subsection{阅读论文时的统一检查框架}
阅读量化论文时，可以先把方法翻译成同一组数学问题：量化哪些张量，采用何种码本，哪些元素共享参数，用什么数据和损失选择编码，是否改变浮点模型的参数化，以及最终使用什么算术内核。然后分别核对理论假设、实验协议和实现约束。这样的阅读方法能把看似繁多的算法归到网格设计、误差补偿、分布变换、训练适应与执行优化等可比较机制。

参考文献中的整数推理与 QAT 基础适合联系第\ref{sec:quant-integer}、\ref{sec:quant-qat}节阅读；AdaRound 与 GPTQ 对应舍入耦合和重建优化；SmoothQuant、AWQ 与 QuaRot 对应改变坐标表示后再量化；QLoRA 对应冻结量化基座上的训练；KIVI 对应在线缓存。本文不把某篇论文在特定模型上的结果提升为普遍保证，也不以论文名代替格式定义。掌握这些机制后，新的方法可以通过它实际改变了哪个目标、哪类误差或哪条执行路径来理解。

\begin{thebibliography}{99}
\bibitem{jacob2018}
Benoit Jacob 等.
Quantization and Training of Neural Networks for Efficient Integer-Arithmetic-Only Inference.
CVPR, 2018.
\url{https://arxiv.org/abs/1712.05877}.

\bibitem{micikevicius2022}
Paulius Micikevicius 等.
FP8 Formats for Deep Learning.
arXiv:2209.05433, 2022.
\url{https://arxiv.org/abs/2209.05433}.

\bibitem{nagel2020}
Markus Nagel 等.
Up or Down? Adaptive Rounding for Post-Training Quantization.
ICML, 2020.
\url{https://arxiv.org/abs/2004.10568}.

\bibitem{esser2020}
Steven K. Esser 等.
Learned Step Size Quantization.
ICLR, 2020.
\url{https://arxiv.org/abs/1902.08153}.

\bibitem{frantar2023}
Elias Frantar, Saleh Ashkboos, Torsten Hoefler, Dan Alistarh.
GPTQ: Accurate Post-Training Quantization for Generative Pre-trained Transformers.
ICLR, 2023.
\url{https://arxiv.org/abs/2210.17323}.

\bibitem{xiao2023}
Guangxuan Xiao 等.
SmoothQuant: Accurate and Efficient Post-Training Quantization for Large Language Models.
ICML, 2023.
\url{https://arxiv.org/abs/2211.10438}.

\bibitem{lin2024}
Ji Lin 等.
AWQ: Activation-aware Weight Quantization for LLM Compression and Acceleration.
arXiv:2306.00978, 2023.
\url{https://arxiv.org/abs/2306.00978}.

\bibitem{dettmers2022}
Tim Dettmers, Mike Lewis, Younes Belkada, Luke Zettlemoyer.
LLM.int8(): 8-bit Matrix Multiplication for Transformers at Scale.
arXiv:2208.07339, 2022.
\url{https://arxiv.org/abs/2208.07339}.

\bibitem{dettmers2023}
Tim Dettmers, Artidoro Pagnoni, Ari Holtzman, Luke Zettlemoyer.
QLoRA: Efficient Finetuning of Quantized LLMs.
arXiv:2305.14314, 2023.
\url{https://arxiv.org/abs/2305.14314}.

\bibitem{ashkboos2024}
Saleh Ashkboos 等.
QuaRot: Outlier-Free 4-Bit Inference in Rotated LLMs.
arXiv:2404.00456, 2024.
\url{https://arxiv.org/abs/2404.00456}.

\bibitem{liu2024}
Zirui Liu 等.
KIVI: A Tuning-Free Asymmetric 2bit Quantization for KV Cache.
arXiv:2402.02750, 2024.
\url{https://arxiv.org/abs/2402.02750}.
\end{thebibliography}
\end{document}
